\textbf{Figure 6. The MASLD Gene Catalog places known MASLD genes, the 26 genes supported in both domains and clinical drug targets by mechanism, disease liability and cell source.}

Genetic counts, overlap estimates, accessibility context and Unified Gene Catalog classes reflect the allele-aware, positive-semidefinite (PSD) LD-projected eQTL refit of the 90 genes with duplicate variant identifiers (\texttt{EQTL-DEDUP-PSD-20261118}: 84 kept positive, 4 lost [CDK5RAP3, KANSL1-AS1, LRRC37A, THBS3], 2 untestable [CTSS, MPRIP]; 430 primary multi-signal SuSiE genes and 990 union genes; 5/5 unaffected controls pass parity at \textbar $\Delta$PP.H4\textbar{} $\le$ 2.55 $\times$ 10$^{-4}$).


Panels (B) to (G) use the Unified Gene Catalog (\texttt{gene-catalog-unified-20261118}; PSD eQTL refit plus Tier-1 direct-MASLD/PDFF/cirrhosis coding rule at PIP $\ge$ 0.5, and ambient-corrected C3b within-hepatocyte partition); (A) and (H) describe the Catalog schema and access model; (G) trial columns come from the curated randomized trial table. 

\textbf{(A)} Catalog structure. Each gene identity carries a separate eQTL-colocalization call and bulk RNA call, each qualified by testability and provenance, plus one molecular object (protein-coding gene, lncRNA transcript, ambiguous lncRNA locus or untestable nomination; no record is a DNA element) and one proposed discriminating experiment. Protein, chromatin and spatial context keep their assay or program unit and enter no class or route.

\textbf{(B)} Fourteen established MASLD genes in the historical adopted Catalog and under the Unified Catalog rule, which admits a coding variant in a fine-mapped credible set at PIP $\ge$ 0.5 specifically when its top GWAS trait is a Tier-1 direct MASLD, liver-fat (PDFF) or cirrhosis trait. PNPLA3 (PIP 0.71), TM6SF2 (1.00), MTARC1 (0.76), APOE (1.00) and SERPINA1 (0.60) move from unresolved to genetically anchored (together with ADH1B and PNPLA2; GCKR and SAMM50 are already regulatory, so all 9 Tier-1 coding genes are anchored). GPAM is regulatory-anchored with disease association at the anchoring SNP. APOB, CIDEB, MTTP, PSD3 and TRIB1 stay unresolved, and HSD17B13 (coding PIP 0.04) and MBOAT7 stay untested. All 17 non-regulatory Tier-2 liver-enzyme-only coding hits (including HBB, GOT1, ABCB11, EPPK1 and FOXN4) are retained outside genetically anchored and both-supported classes.

\textbf{(C)} All 26 genes supported in both domains in the Unified Catalog (all 26 are regulatory liver-enzyme colocalizations: ALT 12, GGT 8, AST 6; Tier-2 coding hits ABCB11, EPPK1 and FOXN4 are excluded). Columns: genetic mechanism; anchor trait; disease liability at the anchoring SNP oriented to the enzyme-raising allele (detected 4: IFI27L2, SPP1, FADS2, IFI30; weak or indeterminate 8; opposite direction 2: HKDC1, KRT8; no association 10; not classifiable, East Asian anchor, 2); bulk RNA, disease versus control (up 21, down 5); and C3b ambient-corrected stage partition of SteatoSITE-replicated stage genes (Figure 3H; inside hepatocytes 3: CFHR4, CFHR5, PCOLCE2; not separable from cell mix 9; unresolved 5; not a replicated stage gene 9). The class was named ``concordant'' in earlier releases; it records a call in both domains and does not claim directional agreement.

\textbf{(D)} MVP NAFLD (filled) and MVP cirrhosis (open) z at the anchoring SNP for the 26 genes in (C), in the anchor's ancestry, oriented to the allele that raises the anchor liver-enzyme trait. Dashed lines, \textbar z\textbar{} = 3.29. No point for East Asian anchors (CDH6, MLIP), which have no same-SNP, same-ancestry European MVP disease lookup.

\textbf{(E)} Unified Catalog fields for 20 genes encoding targets of approved or clinical-stage MASH drugs, grouped by drug mechanism. THRB (regulatory, UK Biobank GGT, PP.H4 \textgreater{} 0.999; no disease association at the anchor), ACACB (regulatory, BBJ GGT, PP.H4 0.75; East Asian anchor, not classifiable) and PNPLA3 (Tier-1 coding, PIP 0.71) are genetically anchored. GIPR, FGF21, FASN, SCD and CCR5 are higher in disease without genetic support. GLP1R is untested (not measured in bulk RNA). Under the ambient-corrected C3b partition, GCGR is classified as not separable from cell mix and has no bulk disease-versus-control call.

\textbf{(F)} Thirty-six genes in three therapeutic axes. Bars, number of domains with a call (bulk RNA, genetics under the Unified Catalog rule, unified MASLD-only liver protein BH q \textless{} 0.05). Transcriptomics, log2FC where BH q \textless{} 0.05, white otherwise: disease versus control (774 genetic-identity clusters), NAS 5--8 versus NAS 0 (46 versus 167), F3 versus F0 (103 versus 69) and F4 versus F0 (103 versus 25). Genetics, SuSiE PP.H4 (gray, no SuSiE result) and Tier-1 coding PIP (white, no coding variant in a credible set). Protein, PXD051911 unified MASLD-only log2FC where BH q \textless{} 0.05; gray, not measured. Right column, clinical status. KRT8 is the only gene with calls in all three domains.

\textbf{(G)} Randomized trial readouts by drug mechanism beside Unified Catalog genetic evidence. Each row is one of 28 drug mechanisms with a randomized MASLD or MASH trial in the curated trial table (2026-11-18 remediation correcting aldafermin to \emph{FGF19}/\emph{FGFR4}/\emph{KLB} and MSDC-0602K to insulin sensitizer [MPC modulator]). Rows are ordered alphabetically by mechanism name; the order carries no priority, and no column is combined into a score. Target genes: one tile per target gene listed for the mechanism (53 target-gene tiles across 43 distinct genes, plus 1 not-applicable tile for vitamin E = 54 tiles), filled by the genetic evidence state in the Unified Gene Catalog (\texttt{gene-catalog-unified-20261118}; PSD eQTL refit plus Tier-1 direct-MASLD/PDFF/cirrhosis coding rule at PIP $\ge$ 0.5). Supported means a regulatory colocalization (tier-1/2 SuSiE) or a Tier-1 coding variant in a credible set at PIP $\ge$ 0.5; indeterminate means testable without a call; untestable means the Catalog could not test the gene; not applicable means the mechanism has no target gene. Counts: supported, regulatory (COLOC) 2 (\emph{ACACB}, SuSiE PP.H4 0.755, BBJ GGT; \emph{THRB}, SuSiE PP.H4 1.000, UKBB GGT); supported, coding (Tier-1, PIP $\ge$ 0.5) 1 (\emph{PNPLA3}, coding PIP 0.71, MVP NAFLD AMR); indeterminate 49 (including \emph{FGF19} and \emph{KLB}); untestable 1 (\emph{FGFR4}); not applicable 1 (vitamin E). In the historical November 6 release \texttt{gene-catalog-eligibility-20261106}, which has no coding rule, \emph{PNPLA3} is indeterminate; no other tile differs. Trial phase: highest randomized trial phase among all randomized trials of the mechanism in the trial table (phase 2: 9 mechanisms; 2b: 10; 3: 7; 4: 2). The table does not record approval status, so phase is not approval. Endpoints: result versus control in the one representative trial per mechanism and endpoint (largest latest-phase randomized trial) for liver fat (MRI-PDFF, MR spectroscopy or an unnamed MRI method, as reported), ALT, MASH resolution and fibrosis improvement. Improved, the trial's own prespecified test (primary or key secondary, with its multiplicity correction) favours the drug; improved, unadjusted, a nominal test favours the drug outside a multiplicity-controlled hierarchy in the source read; not improved, that test did not favour the drug (no cell was coded worsened); no test reported, numbers without a between-group test; dashed outline, no randomized readout in the trial table. Counts over 112 cells: improved 18; improved, unadjusted 19; not improved 25; no test reported 4; no randomized readout 46. Dot, the representative trial row has at least one field marked not verified (11 cells). Trial results are source-owned public data, not findings of this paper; effects are in each trial's own units in the source table. Not plotted: gene-level support at the target genes (MAGMA gene \emph{p} \textless{} 0.05/43 for liver fat and cirrhosis, or any variant at \emph{p} \textless{} 5 $\times$ 10$^{-8}$ within the gene $\pm$50 kb, the only rule for ALT) reaches 3 mechanisms for liver fat (\emph{DGAT2, KHK, PNPLA3}), 1 for cirrhosis (\emph{PNPLA3}), 10 for Pan-UKBB ALT and 6 for MVP ALT, of 27 with a target (\emph{FGF19} has no support on any trait); the matched comparisons (liver fat by liver-fat endpoint, 14 mechanisms, Fisher \emph{p} = 1.0; cirrhosis by fibrosis improvement, 20, \emph{p} = 1.0; ALT by ALT response, 11, \emph{p} = 0.545) give Holm \emph{p} = 1.0 for all three, and none is informative. Mechanisms sharing a target gene (for example \emph{GLP1R} in four incretin rows) are not independent, and absence of genetic support at a target gene is not evidence against the target.

\textbf{(H)} Ways to use the Catalog: view the worked entries, inspect each molecular object with its uncertainty, read source references, and download the gene calls with the molecular-object records. All views draw on the same inputs and give no combined score or rank. The Catalog files are available at \texttt{CATALOG\_DEPOSIT\_ACCESSION\_TBD} (see Data availability).
